\chapter{Introduction}

In many modern applications, systems continuously record observations, measurements, or events over time, and the recorded data naturally form a time-series. Nearly every scientific discipline involves analyzing tasks using time-series data, ranging from the natural and physical sciences to the life and medical sciences, as well as economics, engineering, and the social sciences \cite{boniol2024ts-ad-review}. Thus, time-series anomaly detection is crucial for identifying deviations from normal behavior, which may signal fraud in transaction systems, sensor failures in industrial environments, or urgent changes that require intervention in health care monitoring.

One particular example addresses the difficulty of detecting outliers in large-scale IoT environments, where multiple sensors generate noisy, high-dimensional, and highly correlated data. While traditional threshold-based or statistical methods can help detect anomalies in the IoT industrial environment, they lack the capacity to handle nonlinear, multivariate, and evolving time-series. Consequently, the focus of recent research has shifted towards machine learning and deep learning frameworks that can automatically model complex temporal patterns in sensor data \cite{li2019madgan} \cite{yin2020iot-one-ts-ad} \cite{tuli2022iot-two-ts-ad}.

Another example worth mentioning is a study that addresses the detection of epileptic seizures from EEG recordings, crucial in neuroscience and practically challenging as brain signals are very noisy, nonstationary, and patient-specific \cite{wang2024seizure-ts-ad}. Manual inspection of such deviations is labor-intensive, while conventional machine-learning methods often fail to generalize across patients. Therefore, advancing research on robust anomaly detection methods has the potential to contribute to the broader progress of healthcare and neuroscience.

While the importance of anomaly detection in time-series is indisputable, we should concern ourselves with what we define as an anomaly. Considering that every piece of information, whether digital or in natural form, is generated by "normal" processes, the unusual behavior of these processes creates abnormalities, discordants, deviants, or anomalies \cite{aggarwal2016an-introduction-to-outlier-analysis}. In the context of time-series, a point anomaly is an unusual value related to the global or local distribution of normal values, while a collective anomaly represents a sequence of anomalous points deviating from the expected behavior when observed together. This rather philosophical matter is tackled in detail in Section \ref{sec:ad-for-ts}.

Traditional anomaly detection methods, such as statistical prediction-based methods \cite{pena2013arima-ad} in which an outlier is identified if it falls outside an expected range of values, do not consider contextual anomalies, which are anomalies falling inside the range of those expected values. Other classical methods \cite{scholkopf1999oc-svm} \cite{liu2008iforest} cannot handle the nonlinearity and the dynamic nature of multivariate time-series data. Alternatively, deep learning methods were explored \cite{malhotra2015lstm-ad} \cite{li2019madgan} \cite{an2015vae-ae-ad} for this task, many of which fall under the prediction-based approaches: e.g., Long-Short Term Memory, Autoencoders, or Generative Adversarial Networks, which have demonstrated strong robustness in detecting anomalies.

Generative Adversarial Networks (GANs) represent a deep learning architecture composed of two competing neural networks: a generator, which attempts to produce synthetic samples that resemble normal data, and a discriminator, which learns to distinguish between real and generated samples. These networks are attractive because they can model complex nonlinear dependencies and high-dimensional temporal relationships without relying on explicit assumptions about the data distribution \cite{goodfellow2014gan}. In the context of time-series, the assumptions are about stationarity, the trend, and the seasonal components. Additionally, for anomaly detection, the discriminator provides a score that can be directly used to assess the normality of the data, which is very convenient.

Even though these methods were built previously \cite{li2019madgan} \cite{geiger2020tadgan} \cite{bashar2020tanogan}, the effectiveness of the anomaly scores obtained by GANs in detecting point, contextual, and collective anomalies remains an open research direction, alongside the broader challenge of stabilizing GAN training to fully exploit the representational power of these models for time-series anomaly detection. To address these challenges, in this work, we evaluate two adversarial architectures, MAD-GAN and TadGAN, against a deterministic baseline composed of ARIMA, a standard LSTM, and an LSTM autoencoder. Furthermore, we investigate the performance of various anomaly scoring methods across different anomaly types using three benchmark datasets: NASA space telemetry, Yahoo S5, and the Numenta Anomaly Benchmark (NAB).

\vspace{1em}

% ----------------------------------------------
\textbf{Research Questions} \-\hspace{0.15em} The thesis is focusing on the following research questions:

\begin{enumerate}
    \item How accurately can GAN-based networks reproduce normal time-series patterns compared to other predictive methods, and how does their reconstruction quality correlate with anomaly detection performance? 
    
    \item How well do GAN-based anomaly detection methods perform on collective anomalies with global and contextual deviations, under mean-shift and variance-shift conditions?

    \item How much can a CycleGAN-based approach improve anomaly detection performance compared to a standard GAN in the context of time-series?
\end{enumerate}
% ----------------------------------------------


% ----------------------------------------------
\textbf{Contributions} \-\hspace{0.15em} The contributions of this thesis are:

\begin{enumerate}
    \item \textbf{Cycle Loss Influence:} A fair comparison between standard GAN and cycle-GAN, both trained using Wasserstein loss and evaluated with the same methodology, to determine the importance of learning a circular mapping of time-series in the context of anomaly detection.

    \item \textbf{Prediction Error Evaluation:} An analysis of two different methods to compute anomaly scores using the reconstruction and discriminator scoring capabilities of the GAN, based on a literature review of existing methods.
    
    \item \textbf{Open-Source Benchmark:} Implemented a PyTorch framework that contains four deep learning, predictive-based methods for time-series anomaly detection: LSTM forecasting network, LSTM Autoencoder, adapted MAD-GAN, and Tad-GAN. 
\end{enumerate}
% ----------------------------------------------


% ----------------------------------------------
\textbf{Thesis Outline} \-\hspace{0.15em} The upcoming chapters are structured as follows. Chapter \ref{chap:02-preliminary} provides the base knowledge needed to understand anomaly detection for time-series and GANs. Chapter \ref{chap:03-related-work} describes existing outlier detection methods based on GANs and cycle-GANs. Chapter \ref{chap:04-time-series-anomaly-detection-with-gans} depicts the design, architecture, and training procedure used for the GAN models, as well as how to compute anomaly scores. Chapter \ref{chap:05-evaluation} details the experimental setup, including the datasets used for evaluation and the performance of the baseline methods. Chapter \ref{chap:06-conclusion} completes the thesis by drawing a conclusion and outlining further research. 
% ----------------------------------------------